\chapter{ฟิสิกส์การคายประจุและเคมีพลาสมา}
\label{chap:ch02}

% =========================================================================
% แผนบริหารการสอนประจำบทที่ 2 (Modular Lesson Plan Boxes - 2-Page Standard)
% =========================================================================

\section*{แผนบริหารการสอนประจำบทที่ 2}
\addcontentsline{toc}{section}{แผนบริหารการสอนประจำบทที่ 2}

% --- Page 1 of Lesson Plan ---
\begin{planbox}{หัวข้อเนื้อหาประจำบท}
\begin{enumerate}
    \item ฟิสิกส์การคายประจุในก๊าซและกลไกการพังทลายทางไฟฟ้าของทาวน์เซนด์
    \item กฎของพาสเชนและกลไกสตรีมเมอร์ที่ความดันบรรยากาศ
    \item กระบวนการชน จลนพลศาสตร์ปฏิกิริยา และฟังก์ชันการกระจายพลังงานของอิเล็กตรอน
    \item เคมีพลาสมาและเส้นทางการสังเคราะห์อนุมูลอิสระของออกซิเจนและไนโตรเจน
    \item อันตรกิริยาบริเวณรอยต่อระหว่างพลาสมาและของเหลว
    \item การประยุกต์ใช้โมเดลโครงข่ายประสาทเทียมฟิสิกส์เหนี่ยวนำในการทำนายเคมีพลาสมา
\end{enumerate}
\end{planbox}

\begin{planbox}{วัตถุประสงค์เชิงพฤติกรรม}
\begin{enumerate}
    \item อธิบายกลไกการเกิดขบวนถล่มทาวน์เซนด์และคำนวณเงื่อนไขการพังทลายทางไฟฟ้าได้อย่างถูกต้อง
    \item วิเคราะห์เส้นโค้งพาสเชนและระบุสภาวะความต่างศักย์พังทลายต่ำสุดสำหรับก๊าซชนิดต่างๆ ได้
    \item อธิบายกลไกสตรีมเมอร์และคำนวณเกณฑ์เรเทอร์-มีกที่ความดันบรรยากาศได้อย่างแม่นยำ
    \item เขียนสมการเคมีและวิเคราะห์เส้นทางการเกิดอนุภาค RONS ชนิดหลักได้อย่างครบถ้วน
    \item ประยุกต์ใช้โมเดล Deep Learning แบบ PINN ในการจำลองจลนพลศาสตร์เคมีพลาสมาแบบเรียลไทม์ได้
\end{enumerate}
\end{planbox}

\newpage

% --- Page 2 of Lesson Plan ---
\begin{planbox}{วิธีสอนและกิจกรรมการเรียนการสอน}
\begin{enumerate}
    \item บรรยายเชิงทฤษฎีควบคู่กับการสาธิตแบบจำลองการจำลองสตรีมเมอร์ 2 มิติ
    \item คำนวณความต่างศักย์พังทลายตามกฎของพาสเชนในตัวอย่างการออกแบบหัวพ่นพลาสมาเจ็ท
    \item วิเคราะห์เมทริกซ์ปฏิกิริยาเคมี RONS และการตรวจวัดทางสเปกโทรสโกปี
    \item ฝึกปฏิบัติการเขียนโค้ดโครงข่ายประสาทเทียม PINN สำหรับแก้สมการเชิงอนุพันธ์ของเคมีพลาสมา
    \item ทำงานกลุ่มย่อยเพื่อวิเคราะห์ความสัมพันธ์ระหว่างค่า pH, ORP และความเข้มข้นของไนเตรตในน้ำ PAW
\end{enumerate}
\end{planbox}

\begin{planbox}{สื่อการเรียนการสอน}
\begin{enumerate}
    \item เอกสารคำสอนและโปรแกรมแบบจำลอง Boltzmann Solver (BOLSIG+)
    \item โค้ดตัวอย่างภาษาไพทอน (PyTorch) สำหรับสถาปัตยกรรม Physics-Informed Neural Network
    \item กราฟฟิกเวกเตอร์แสดงวิวัฒนาการของหัวสตรีมเมอร์และเส้นทางปฏิกิริยาเคมีพลาสมา
    \item ชุดข้อมูลการวัดค่าทางเคมีและฟิสิกส์จริงจากห้องปฏิบัติการพลาสมาเกษตรดิจิทัล
\end{enumerate}
\end{planbox}

\begin{planbox}{การวัดและประเมินผล}
\begin{enumerate}
    \item ประเมินจากการมีส่วนร่วมในการแก้โจทย์ฟิสิกส์การคายประจุในชั้นเรียน ร้อยละ 10
    \item ประเมินจากความถูกต้องของรายงานการจำลองเคมี RONS ด้วยโมเดล PINN ร้อยละ 20
    \item ประเมินจากการทำแบบฝึกหัดทบทวนท้ายบท ร้อยละ 10
    \item ประเมินจากการสอบข้อเขียนประจำภาคการศึกษา ร้อยละ 60
\end{enumerate}
\end{planbox}

\clearpage

% =========================================================================
% เนื้อหาบทเรียน (เริ่มต้นหน้า 3 อย่างเป็นระเบียบ)
% =========================================================================

\section{ฟิสิกส์การคายประจุในก๊าซ}

การคายประจุในก๊าซ (Gas Discharge) คือกระบวนการที่กระแสไฟฟ้าไหลผ่านตัวกลางก๊าซอันเนื่องมาจากการแตกตัวเป็นไอออนของอะตอมหรือโมเลกุลภายใต้อิทธิพลของสนามไฟฟ้าภายนอก การทำความเข้าใจฟิสิกส์การคายประจุเป็นพื้นฐานสำคัญที่สุดในการออกแบบแหล่งกำเนิดพลาสมาให้ทำงานได้อย่างมีเสถียรภาพ ประหยัดพลังงาน และสามารถควบคุมการผลิตอนุมูลอิสระ RONS ได้ตามความต้องการทางชีวภาพ

\subsection{การพังทลายทางไฟฟ้าและขบวนถล่มของทาวน์เซนด์}

ทฤษฎีการคายประจุของทาวน์เซนด์ (Townsend Breakdown Theory) อธิบายกระบวนการเปลี่ยนผ่านจากก๊าซที่เป็นฉนวนไฟฟ้าไปสู่สถานะนำไฟฟ้า โดยพิจารณาอิเล็กตรอนตั้งต้น (Seed electron) ที่เกิดขึ้นจากรังสีคอสมิกหรือรังสีพื้นหลัง เมื่อได้รับแรงเร่งจากสนามไฟฟ้าภายนอก อิเล็กตรอนจะพุ่งเข้าชนโมเลกุลของก๊าซอย่างต่อเนื่อง หากพลังงานจลน์ของอิเล็กตรอนสูงกว่าพลังงานแตกตัวเป็นไอออน (Ionization potential) จะทำให้อิเล็กตรอนวงนอกหลุดออก เกิดเป็นอิเล็กตรอนตัวใหม่และไอออนบวกหนึ่งคู่

กระบวนการนี้จะทวีคูณจำนวนอย่างรวดเร็วในลักษณะลูกโซ่ เรียกว่า \textbf{ขบวนถล่มอิเล็กตรอน (Electron Avalanche)} อัตราการเพิ่มขึ้นของจำนวนอิเล็กตรอนต่อหน่วยระยะทางในทิศทางของสนามไฟฟ้า กำหนดโดย \textbf{สัมประสิทธิ์การแตกตัวขั้นแรกของทาวน์เซนด์ (Townsend First Ionization Coefficient, $\alpha$)}
\begin{equation}
    dn_e = \alpha \, n_e \, dx
\end{equation}
เมื่ออินทิเกรตตามระยะห่างระหว่างขั้วไฟฟ้า $d$ จะได้จำนวนอิเล็กตรอนที่เคลื่อนที่ไปถึงขั้วบวก (Anode) คือ
\begin{equation}
    n_e(d) = n_0 \exp(\alpha d)
\end{equation}
โดยที่ $n_0$ คือจำนวนอิเล็กตรอนตั้งต้นที่หลุดออกจากขั้วลบ (Cathode)

ในเวลาเดียวกัน ไอออนบวกที่เกิดขึ้นจะเคลื่อนที่สวนทิศทางกลับไปยังขั้วลบ เมื่อไอออนบวกพุ่งชนผิวขั้วลบ จะมีโอกาสปลดปล่อยอิเล็กตรอนทุติยภูมิ (Secondary electrons) ออกมา เรียกว่า \textbf{สัมประสิทธิ์การปล่อยอิเล็กตรอนทุติยภูมิของทาวน์เซนด์ (Townsend Secondary Ionization Coefficient, $\gamma$)} 

เงื่อนไขที่การคายประจุจะสามารถดำรงอยู่ได้ด้วยตนเอง (Self-sustaining Discharge Criterion) โดยไม่ต้องพึ่งพาแหล่งกำเนิดอิเล็กตรอนภายนอกอีกต่อไป คือผลคูณของจำนวนไอออนที่เกิดขึ้นกับความน่าจะเป็นในการปลดปล่อยอิเล็กตรอนทุติยภูมิต้องมีค่าอย่างน้อยเท่ากับ 1
\begin{equation}
    \gamma \left( \exp(\alpha d) - 1 \right) = 1
    \label{eq:townsend_criterion}
\end{equation}
ความสัมพันธ์ในสมการที่ \eqref{eq:townsend_criterion} คือเกณฑ์การพังทลายทางไฟฟ้าแบบทาวน์เซนด์ (Townsend Breakdown Criterion)

\subsection{กฎของพาสเชน}

เมื่อนำสัมประสิทธิ์ $\alpha$ ซึ่งเป็นฟังก์ชันของความเข้มสนามไฟฟ้าต่อความหนาแน่นก๊าซ ($E/N$ หรือ $E/p$) ในรูป $\alpha/p = A \exp(-B p / E)$ แทนลงในเกณฑ์ของทาวน์เซนด์ จะได้สมการความสัมพันธ์ระหว่างความต่างศักย์พังทลาย (Breakdown Voltage, $V_B$) กับผลคูณของความดันก๊าซและระยะห่างระหว่างขั้วไฟฟ้า ($p \cdot d$) ซึ่งเรียกว่า \textbf{กฎของพาสเชน (Paschen's Law)}
\begin{equation}
    V_B = \frac{B \cdot (p d)}{\ln(A \cdot p d) - \ln\left[\ln\left(1 + \frac{1}{\gamma}\right)\right]}
    \label{eq:paschen_law}
\end{equation}
โดยที่ $A$ และ $B$ เป็นค่าคงที่เฉพาะของก๊าซแต่ละชนิด

\begin{figure}[H]
\centering
\begin{tikzpicture}[scale=0.92]
    % Axes
    \draw[->, >=Stealth, line width=1.1pt, plasmaNavy] (0,0) -- (7.5,0) node[right, font=\small\bfseries] {$p \cdot d\quad (\text{Torr}\cdot\text{cm})$};
    \draw[->, >=Stealth, line width=1.1pt, plasmaNavy] (0,0) -- (0,5.2) node[above, font=\small\bfseries] {$V_B\quad (\text{V})$};

    % Grid
    \draw[gray!20, dashed, step=1cm] (0,0) grid (7,4.8);

    % Paschen curves
    % Air Curve
    \draw[line width=1.5pt, red!80!black] 
        (0.6, 4.6) .. controls (1.2, 1.8) and (1.6, 1.3) .. (2.0, 1.35)
        .. controls (2.8, 1.45) and (4.5, 2.7) .. (6.8, 4.5)
        node[right, font=\small\bfseries, red!80!black] {อากาศ (Air)};

    % Argon Curve
    \draw[line width=1.4pt, plasmaCyan!90!black] 
        (0.8, 4.4) .. controls (1.5, 1.3) and (2.0, 0.85) .. (2.5, 0.9)
        .. controls (3.5, 1.1) and (5.0, 2.0) .. (6.8, 3.4)
        node[right, font=\small\bfseries, plasmaCyan!90!black] {อาร์กอน (Ar)};

    % Helium Curve
    \draw[line width=1.4pt, plasmaViolet] 
        (1.0, 4.2) .. controls (1.8, 0.9) and (2.5, 0.55) .. (3.2, 0.6)
        .. controls (4.2, 0.8) and (5.5, 1.4) .. (6.8, 2.2)
        node[right, font=\small\bfseries, plasmaViolet] {ฮีเลียม (He)};

    % Paschen minimum marker for Air
    \draw[dashed, red!80!black, line width=0.8pt] (2.0, 0) -- (2.0, 1.35) -- (0, 1.35);
    \fill[red!80!black] (2.0, 1.35) circle (2.5pt);
    \node[below, font=\footnotesize, red!80!black] at (2.0, -0.05) {$(pd)_{\min}$};
    \node[left, font=\footnotesize, red!80!black] at (0, 1.35) {$V_{B,\min} \approx 327\text{ V}$};

    % Region explanations
    \node[align=center, font=\footnotesize, gray!80!black] at (1.1, 3.2) {อิเล็กตรอนชนก๊าซน้อย\\(เส้นทางอิสระเฉลี่ยยาว)};
    \node[align=center, font=\footnotesize, gray!80!black] at (5.2, 3.8) {อิเล็กตรอนชนถี่เกินไป\\(สูญเสียพลังงานจลน์)};
\end{tikzpicture}
\caption{เส้นโค้งพาสเชนเปรียบเทียบความต่างศักย์พังทลายทางไฟฟ้าระหว่างอากาศ อาร์กอน และฮีเลียม}
\label{fig:paschen_curve}
\end{figure}

ลักษณะสำคัญของเส้นโค้งพาสเชนคือ การมี \textbf{จุดต่ำสุดของพาสเชน (Paschen Minimum)}
\begin{itemize}
    \item \textbf{ทางด้านซ้ายของจุดต่ำสุด ($p \cdot d < (pd)_{\min}$):} ก๊าซเบาบางหรือระยะห่างแคบเกินไป ทำให้อิเล็กตรอนเคลื่อนที่ผ่านช่องว่างโดยแทบไม่ชนโมเลกุลก๊าซ จึงต้องใช้ศักย์ไฟฟ้าสูงมากในการเหนี่ยวนำให้เกิดการแตกตัว
    \item \textbf{ทางด้านขวาของจุดต่ำสุด ($p \cdot d > (pd)_{\min}$):} ก๊าซหนาแน่นหรือระยะห่างกว้าง ทำให้อิเล็กตรอนชนกับโมเลกุลก๊าซบ่อยครั้งเกินไปจนสูญเสียพลังงานจลน์สะสม จึงต้องใช้ศักย์ไฟฟ้าสูงขึ้นเพื่อชดเชยพลังงานที่สูญเสียไปจากการชน
\end{itemize}

ในอากาศที่ความดันบรรยากาศ ($p \approx 760\text{ Torr}$) หากระยะห่างขั้วไฟฟ้า $d = 1\text{ mm}$ ผลคูณ $pd = 76\text{ Torr}\cdot\text{cm}$ ซึ่งอยู่ทางขวาของจุดต่ำสุดมาก ทำให้ต้องใช้ความต่างศักย์พังทลายสูงถึงประมาณ $3 - 4\text{ kV}$ ในการเริ่มต้นการคายประจุ

\subsection{กลไกสตรีมเมอร์ที่ความดันบรรยากาศ}

เมื่อผลคูณ $pd$ มีค่าสูง (เช่น ที่ความดันบรรยากาศ 1 atm) ทฤษฎีของทาวน์เซนด์จะไม่สามารถอธิบายความเร็วในการลุกลามของการคายประจุได้อย่างสมบูรณ์ เนื่องจากพบว่าการพังทลายเกิดขึ้นในเวลาเพียงไม่กี่นาโนวินาที ซึ่งเร็วกว่าความเร็วการเคลื่อนที่ของไอออนบวกกลับไปยังขั้วลบนับพันเท่า

ปรากฏการณ์นี้อธิบายด้วย \textbf{ทฤษฎีสตรีมเมอร์ (Streamer Theory)} ของมีกและเรเทอร์ (Meek and Raether) เมื่อขบวนถล่มอิเล็กตรอนเติบโตจนมีจำนวนประจุสะสมที่หัวขบวนเกินกว่า $10^8$ ตัว สนามไฟฟ้าที่เกิดจากประจุอวกาศเฉพาะที่ (Space Charge Field, $E_{sc}$) จะมีขนาดเทียบเท่ากับสนามไฟฟ้าภายนอก ($E_{sc} \approx E_0$) ส่งผลให้เกิดการแผ่โฟตอนพลังงานสูง (Photoionization) ไปแตกตัวก๊าซด้านหน้า เกิดเป็นลำไอออนนำไฟฟ้าความเร็วสูงพุ่งเชื่อมระหว่างขั้วไฟฟ้า เรียกว่า \textbf{เกณฑ์เรเทอร์-มีก (Raether-Meek Criterion)}
\begin{equation}
    \alpha \cdot d \approx 18 - 20 \quad \text{หรือ} \quad N_e \ge 10^8
\end{equation}
กลไกสตรีมเมอร์เป็นหัวใจสำคัญในการทำความเข้าใจการคายประจุแบบกั้นด้วยไดอิเล็กทริก (DBD) และพลาสมาเจ็ทความดันบรรยากาศ (APPJ)

\section{กระบวนการชนและจลนพลศาสตร์ปฏิกิริยา}

ปฏิกิริยาเคมีกายภาพทั้งหมดในพลาสมาเย็นเริ่มต้นจากการชนระหว่างอิเล็กตรอนพลังงานสูงกับอะตอมหรือโมเลกุลของก๊าซพื้นฐาน

\subsection{การชนแบบยืดหยุ่นและการชนแบบไม่ยืดหยุ่น}

\begin{enumerate}
    \item \textbf{การชนแบบยืดหยุ่น (Elastic Collisions)} เป็นการชนที่ไม่มีการเปลี่ยนรูปพลังงานไปเป็นพลังงานกระตุ้นภายในอะตอม พลังงานจลน์รวมยังคงอนุรักษ์ เนื่องจากอิเล็กตรอนมีมวลน้อยกว่าโมเลกุลก๊าซมาก สัดส่วนการถ่ายโอนพลังงานในการชนแต่ละครั้งมีค่าน้อยมาก ($\Delta E / E \approx 2 m_e / M \sim 10^{-5}$) ทำให้อิเล็กตรอนสามารถสะสมพลังงานจลน์จนมีอุณหภูมิสูงได้โดยไม่ทำให้อุณหภูมิรวมของก๊าซร้อนขึ้น
    \item \textbf{การชนแบบไม่ยืดหยุ่น (Inelastic Collisions)} อิเล็กตรอนถ่ายโอนพลังงานจลน์ไปเปลี่ยนสถานะทางควอนตัมของอะตอมหรือโมเลกุล ได้แก่
    \begin{itemize}
        \item \textbf{การกระตุ้นระดับพลังงาน (Excitation):} $\ce{e^- + N2 -> e^- + N2^*}$
        \item \textbf{การแตกตัวเป็นไอออนโดยตรง (Direct Ionization):} $\ce{e^- + O2 -> 2e^- + O2^+}$
        \item \textbf{การแตกพันธะโมเลกุล (Dissociation):} $\ce{e^- + O2 -> e^- + O(^3P) + O(^1D)}$
        \item \textbf{การจับยึดอิเล็กตรอน (Electron Attachment):} $\ce{e^- + O2 -> O^- + O}$
    \end{itemize}
\end{enumerate}

\subsection{ฟังก์ชันการกระจายพลังงานของอิเล็กตรอน}

อัตราการเกิดปฏิกิริยาเคมีพลาสมา ($k_j$) ถูกควบคุมโดย \textbf{ฟังก์ชันการกระจายพลังงานของอิเล็กตรอน (Electron Energy Distribution Function หรือ EEDF, $f(\varepsilon)$)}
\begin{equation}
    k_j = \sqrt{\frac{2}{m_e}} \int_0^\infty \varepsilon \, \sigma_j(\varepsilon) \, f(\varepsilon) \, d\varepsilon
\end{equation}
โดยที่ $\sigma_j(\varepsilon)$ คือภาคตัดขวางการชนของปฏิกิริยา (Cross-section) ในพลาสมาความดันบรรยากาศ ฟังก์ชัน EEDF มักเบี่ยงเบนออกจากการกระจายตัวแบบแมกซ์เวลล์ (Non-Maxwellian) โดยมีลักษณะหางการกระจายตัวแบบดรอยเวสเตน (Druyvesteyn-like distribution) ซึ่งมีอิเล็กตรอนพลังงานสูง ($> 10\text{ eV}$) ในสัดส่วนที่เพียงพอต่อการแตกพันธะโมเลกุลที่แข็งแกร่งของ \ce{N2} และ \ce{O2}

\section{เคมีพลาสมาและการสังเคราะห์อนุมูลอิสระ RONS}

หัวใจของการประยุกต์ใช้พลาสมาในงานเกษตรและสิ่งแวดล้อมคือการควบคุมวิถีการสังเคราะห์กลุ่มอนุภาคที่ว่องไวทางเคมีของออกซิเจนและไนโตรเจน (Reactive Oxygen and Nitrogen Species หรือ RONS)

\subsection{เส้นทางการสังเคราะห์อนุภาค ROS}

ออกซิเจนในบรรยากาศเมื่อทำปฏิกิริยากับอิเล็กตรอนพลังงานสูงและไอน้ำ จะก่อให้เกิดอนุมูลอิสระและโมเลกุลออกซิไดซ์สำคัญ ดังสมการ
\begin{align}
    \ce{e^- + O2} &\rightarrow \ce{e^- + O(^3P) + O(^1D)} \quad (\text{การแตกโมเลกุลออกซิเจน}) \label{eq:o_dissoc}\\
    \ce{O(^3P) + O2 + M} &\rightarrow \ce{O3 + M} \quad (\text{การเกิดโอโซน}) \label{eq:o3_form}\\
    \ce{e^- + H2O} &\rightarrow \ce{e^- + \cdot OH + H} \quad (\text{การเกิดอนุมูลไฮดรอกซิล}) \label{eq:oh_form}\\
    \ce{\cdot OH + \cdot OH + M} &\rightarrow \ce{H2O2 + M} \quad (\text{การรวมตัวเป็นไฮโดรเจนเปอร์ออกไซด์}) \label{eq:h2o2_form}
\end{align}
อนุมูลไฮดรอกซิล ($\cdot\ce{OH}$) มีศักย์ไฟฟ้าออกซิเดชันสูงที่สุด ($+2.80\text{ V}$) และมีอายุขัยสั้นมากในระดับไมโครวินาที ทำหน้าที่เป็นสารทำลายผนังเซลล์เชื้อโรคขั้นรุนแรง ในขณะที่โอโซน (\ce{O3}) และไฮโดรเจนเปอร์ออกไซด์ (\ce{H2O2}) มีอายุขัยยาวนานกว่า สามารถละลายลงสู่น้ำเพื่อทำหน้าที่เป็นสารฆ่าเชื้อในระยะยาว

\subsection{เส้นทางการสังเคราะห์อนุภาค RNS}

พลังงานไฟฟ้าในพลาสมาสามารถทำลายพันธะสามของไนโตรเจน ($\ce{N\bond{3}N}$ ซึ่งมีพลังงานพันธะสูงถึง $9.76\text{ eV}$) ทำให้เกิดการตรึงไนโตรเจนจากอากาศสู่สารประกอบว่องไวทางเคมี
\begin{align}
    \ce{e^- + N2} &\rightarrow \ce{e^- + N(^4S) + N(^2D)} \label{eq:n2_dissoc}\\
    \ce{N + O2} &\rightarrow \ce{\cdot NO + O} \quad (\text{กลไกเซลโดวิช}) \label{eq:zeldovich}\\
    \ce{\cdot NO + O3} &\rightarrow \ce{\cdot NO2 + O2} \label{eq:no2_form}\\
    \ce{\cdot NO + \cdot O2^-} &\rightarrow \ce{ONOO^-} \quad (\text{เพอร์ออกซีไนไตรต์}) \label{eq:onoo_form}
\end{align}
เพอร์ออกซีไนไตรต์ (\ce{ONOO^-}) เกิดจากการรวมตัวระหว่างไนตริกออกไซด์และซูเปอร์ออกไซด์ เป็นอนุมูลออกซิไดซ์ทรงพลังที่สามารถซึมผ่านเยื่อหุ้มเซลล์ของแบคทีเรียและเชื้อราเพื่อเหนี่ยวนำให้เกิดกระบวนการตายของเซลล์ (Apoptosis)

\begin{table}[H]
\caption{คุณสมบัติทางเคมีกายภาพและบทบาททางชีววิทยาของอนุภาค RONS หลักในพลาสมา}
\label{tab:rons_summary}
\begin{tabularx}{\textwidth}{p{2.2cm} c c X}
\toprule
\textbf{ชนิดอนุภาค} & \textbf{ศักย์ออกซิเดชัน} & \textbf{ครึ่งชีวิต ($t_{1/2}$)} & \textbf{บทบาทหลักในงานเกษตรและสิ่งแวดล้อม} \\
\midrule
$\cdot\ce{OH}$ & $+2.80\text{ V}$ & $\sim 1\text{ }\mu\text{s}$ & สลายยาฆ่าแมลงตกค้าง, ทำลายผนังเซลล์เชื้อโรค \\
\ce{O3} & $+2.07\text{ V}$ & หลายนาทีถึงชั่วโมง & กำจัดเชื้อราและสปอร์บนผลผลิตสด, สลายสารพิษอะฟลาทอกซิน \\
\ce{H2O2} & $+1.77\text{ V}$ & หลายวัน & สารส่งสัญญาณกระตุ้นการงอกของเมล็ด, เสริมภูมิต้านทานพืช \\
$\cdot\ce{NO}$ & $+1.59\text{ V}$ & ไม่กี่วินาที & สารส่งสัญญาณทำลายระยะพักตัวของเมล็ด, ควบคุมการเปิดปิดปากใบ \\
\ce{ONOO^-} & $+1.40\text{ V}$ & $\sim 10 - 20\text{ ms}$ & ยับยั้งเชื้อราโรครากเน่า, ย่อยสลายสารพิษตกค้าง \\
\ce{NO3^-} & $+0.96\text{ V}$ & เสถียรสูง & แหล่งธาตุอาหารไนโตรเจนหลักสำหรับพืชในน้ำ PAW \\
\bottomrule
\end{tabularx}
\end{table}

\section{อันตรกิริยาบริเวณรอยต่อระหว่างพลาสมาและของเหลว}

เมื่อนำพลาสมามาสัมผัสกับผิวน้ำ เพื่อผลิต \textbf{น้ำกระตุ้นด้วยพลาสมา (Plasma-Activated Water หรือ PAW)} อนุภาค RONS ในสถานะก๊าซจะละลายและถ่ายโอนมวลข้ามรอยต่อ (Interface Transfer) ไปทำปฏิกิริยาเคมีทุติยภูมิในน้ำ
\begin{align}
    \ce{\cdot NO + \cdot NO2 + H2O} &\rightarrow \ce{2 HNO2} \quad (\text{กรดไนตรัส}) \label{eq:hno2}\\
    \ce{2 \cdot NO2 + H2O} &\rightarrow \ce{HNO2 + HNO3} \quad (\text{กรดไนตรัสและกรดไนตริก}) \label{eq:hno3}
\end{align}
ผลลัพธ์ที่เกิดขึ้นทำให้น้ำ PAW มีการเปลี่ยนแปลงคุณสมบัติทางเคมีกายภาพ 4 ประการอย่างชัดเจน
\begin{enumerate}
    \item \textbf{ค่าความเป็นกรด-ด่าง (pH) ลดลง} จากการแตกตัวให้ไฮโดรเนียมไอออน ($\ce{H3O^+}$) ของกรดไนตริกและกรดไนตรัส โดยทั่วไป pH จะลดลงมาอยู่ในช่วง $2.8 - 4.5$
    \item \textbf{ค่าศักย์รีดอกซ์ (Oxidation-Reduction Potential หรือ ORP) สูงขึ้น} จากระดับปกติ $+200\text{ mV}$ พุ่งสูงขึ้นเป็น $+500\text{ mV}$ ถึง $+650\text{ mV}$ บ่งชี้ถึงสภาวะการออกซิเดชันที่รุนแรง
    \item \textbf{ค่าการนำไฟฟ้า (Electrical Conductivity หรือ EC) เพิ่มขึ้น} ตามความเข้มข้นของไอออน \ce{NO3^-}, \ce{NO2^-}, และ \ce{H^+}
    \item \textbf{การสะสมของปุ๋ยไนโตรเจนธรรมชาติ} ไนเตรตที่เกิดขึ้นในน้ำมีความเสถียรสูงและพืชสามารถดูดซึมไปใช้ในการสร้างโปรตีนได้ทันที
\end{enumerate}

\begin{aibox}{การบูรณาการ Deep Learning: Physics-Informed Neural Networks (PINN)}
ระบบเคมีพลาสมาประกอบด้วยปฏิกิริยาเคมีที่เชื่อมโยงกันมากกว่า 50 ชนิด ซึ่งมีอัตราการเกิดปฏิกิริยาแตกต่างกันตั้งแต่ระดับนาโนวินาทีไปจนถึงหลายชั่วโมง ก่อให้เกิดสมการเชิงอนุพันธ์สามัญแบบแข็ง (Stiff Ordinary Differential Equations) ที่การจำลองด้วยวิธีไฟไนต์ดิฟเฟอเรนซ์แบบดั้งเดิมต้องใช้พลังการประมวลผลสูงและทำงานแบบเรียลไทม์ไม่ได้

\textbf{สถาปัตยกรรม PINN สำหรับเคมีพลาสมา:} โครงข่ายประสาทเทียมถูกฝึกสอนให้เรียนรู้ความสัมพันธ์ระหว่างพารามิเตอร์ขาเข้า (แรงดันไฟฟ้า, ความถี่, อัตราไหลของก๊าซ, ความชื้นสัมพัทธ์) กับความเข้มข้นของอนุภาค RONS โดยฟังก์ชันการสูญเสีย (Loss Function) จะถูกผูกมัดด้วยกฎการอนุรักษ์มวลและสมดุลปฏิกิริยาเคมี
\begin{equation}
\begin{aligned}
    \mathcal{L}_{total} &= \mathcal{L}_{data} + \lambda_{kinetics} \cdot \mathcal{L}_{physics} \\
    \mathcal{L}_{physics} &= \sum_{i} \left\| \frac{d[C_i]}{dt} - \left( \mathcal{R}_{prod, i} - \mathcal{R}_{loss, i} \right) \right\|^2
\end{aligned}
\end{equation}
ผลลัพธ์คือ โมเดล AI สามารถทำนายความเข้มข้นของ $\cdot\ce{OH}$ และ \ce{ONOO^-} ในน้ำ PAW ได้แม่นยำกว่าร้อยละ 95 ด้วยเวลาประมวลผลน้อยกว่า 2 มิลลิวินาที ทำให้สามารถติดตั้งบนบอร์ดไมโครคอมพิวเตอร์ขอบ (Edge Device เช่น Raspberry Pi 5 หรือ Jetson Orin Nano) เพื่อควบคุมปริมาณสารออกฤทธิ์ทางชีวภาพในแปลงปลูกจริงได้แบบวงรอบปิด (Closed-loop Control)
\end{aibox}

\begin{examplebox}{ตัวอย่างการคำนวณความต่างศักย์พังทลายตามกฎของพาสเชน}
\textbf{โจทย์} จงคำนวณหาความต่างศักย์พังทลาย ($V_B$) สำหรับระบบการคายประจุในอากาศแห้งที่มีระยะห่างระหว่างขั้วไฟฟ้าระนาบคู่ $d = 2\text{ mm}$ ที่ความดันบรรยากาศปกติ ($p = 760\text{ Torr}$) กำหนดให้ค่าคงที่ของอากาศคือ $A = 15\text{ cm}^{-1}\cdot\text{Torr}^{-1}$, $B = 365\text{ V}\cdot\text{cm}^{-1}\cdot\text{Torr}^{-1}$ และสัมประสิทธิ์การปล่อยอิเล็กตรอนทุติยภูมิ $\gamma = 0.01$

\textbf{การแก้ปัญหา}
\begin{enumerate}
    \item แปลงหน่วยระยะห่าง $d = 2\text{ mm} = 0.2\text{ cm}$
    \item คำนวณผลคูณ $p \cdot d = 760\text{ Torr} \times 0.2\text{ cm} = 152\text{ Torr}\cdot\text{cm}$
    \item คำนวณเทอมตัวส่วนในกฎของพาสเชน
    \begin{align*}
        \ln(A \cdot pd) &= \ln(15 \times 152) = \ln(2280) \approx 7.732 \\
        \ln\left(1 + \frac{1}{\gamma}\right) &= \ln\left(1 + \frac{1}{0.01}\right) = \ln(101) \approx 4.615 \\
        \ln\left[\ln\left(1 + \frac{1}{\gamma}\right)\right] &= \ln(4.615) \approx 1.529
    \end{align*}
    \item แทนค่าลงในสมการกฎของพาสเชน \eqref{eq:paschen_law}
    \begin{equation*}
        V_B = \frac{365 \times 152}{7.732 - 1.529} = \frac{55480}{6.203} \approx 8944\text{ V} \approx 8.94\text{ kV}
    \end{equation*}
\end{enumerate}
\textbf{ความหมายทางกายภาพ} ในการออกแบบระบบหัวพ่นพลาสมาเจ็ทที่ระยะห่าง 2 มิลลิเมตรในอากาศ แหล่งจ่ายไฟฟ้าแรงสูงต้องสามารถจ่ายแรงดันพีคได้อย่างน้อย 8.94 กิโลโวลต์ จึงจะสามารถเริ่มต้นการแตกตัวเป็นไอออนในอากาศได้
\end{examplebox}

\clearpage

% =========================================================================
% คำถามทบทวนท้ายบทที่ 2 (ขึ้นหน้าใหม่อย่างสง่างาม)
% =========================================================================

\section*{คำถามทบทวนท้ายบทที่ 2}
\addcontentsline{toc}{section}{คำถามทบทวนท้ายบทที่ 2}

\begin{enumerate}
    \item จงอธิบายความแตกต่างทางฟิสิกส์ระหว่างกระบวนการขบวนถล่มอิเล็กตรอนของทาวน์เซนด์ (Townsend Avalanche) และการพังทลายแบบสตรีมเมอร์ (Streamer Breakdown) เหตุใดที่ความดันบรรยากาศกลไกสตรีมเมอร์จึงเป็นกลไกหลัก
    \item จากเส้นโค้งพาสเชน เหตุใดการใช้ก๊าซฮีเลียม (He) หรืออาร์กอน (Ar) เป็นก๊าซพาหะในหัวพ่นพลาสมาเจ็ท จึงต้องการแรงดันไฟฟ้าในการจุดพลาสมาต่ำกว่าการใช้อากาศปกติอย่างมีนัยสำคัญ
    \item จงอธิบายกลไกที่ทำให้อิเล็กตรอนในพลาสมาเย็นมีพลังงานจลน์เฉลี่ยสูงถึงหลายอิเล็กตรอนโวลต์ แต่ไม่ทำให้อุณหภูมิของก๊าซร้อนขึ้น โดยอ้างอิงหลักการถ่ายโอนพลังงานในการชนแบบยืดหยุ่น
    \item จงเขียนสมการปฏิกิริยาเคมีที่แสดงการสังเคราะห์เพอร์ออกซีไนไตรต์ (\ce{ONOO^-}) ในพลาสมา และอธิบายว่าเหตุใดอนุภาคนี้จึงมีบทบาทสำคัญอย่างยิ่งในการยับยั้งเชื้อจุลินทรีย์ก่อโรคพืช
    \item การนำสถาปัตยกรรม Physics-Informed Neural Network (PINN) มาประยุกต์ใช้ในการทำนายปฏิกิริยาเคมีพลาสมา มีข้อได้เปรียบเหนือกว่าโมเดลแบบกล่องดำ (Black-box Machine Learning) ทั่วไปอย่างไรในแง่ของความถูกต้องทางวิทยาศาสตร์และความปลอดภัยในการนำไปใช้ในแปลงเกษตรจริง
\end{enumerate}

\clearpage
